\chapter*{Yang's Redshifted Flux of Infinite Light Sources}

\lettrine{L}{ee} Yang had always been captivated by the grandeur of cosmological phenomena. Unlike his colleague Willard Jones, who pondered the flux of stationary light sources, Lee turned his attention to the dynamics of an expanding universe. 

\begin{figure}[H]
\includegraphics[width=1.0\linewidth]{Illustrations/vignette_yang.png}
\caption{Yang turning pages of a book and coffee into thoughts}
\end{figure}

He wondered: What is the total flux of light observed from an infinite grid of light sources if they are all receding according to Hubble's Law?

\section*{The Setup: Hubble’s Law and Redshifted Flux}
In Lee's scenario, each light source moves away from the observer at a speed proportional to its distance \( d \), governed by Hubble’s Law:
\[
v = H_0 d,
\]
where:
\begin{itemize}
    \item \( v \) is the velocity of the light source,
    \item \( H_0 \) is the Hubble constant,
    \item \( d \) is the distance to the light source.
\end{itemize}

The light from these sources undergoes a Doppler shift due to their motion. The observed frequency \( \nu_{\text{obs}} \) of the light is related to the emitted frequency \( \nu_{\text{emit}} \) by:
\[
\nu_{\text{obs}} = \nu_{\text{emit}} \frac{c}{c + v},
\]
where \( c \) is the speed of light.

The observed flux \( F_{\text{obs}} \) from a single light source is affected by this redshift and the inverse square law:
\[
F_{\text{obs}} = \frac{L}{d^2} \cdot \left(\frac{\nu_{\text{obs}}}{\nu_{\text{emit}}}\right)^2,
\]
where \( L \) is the intrinsic luminosity of the source.

% >>>>>> ANNEX D (cut-sheet v2): the total redshifted flux on a unit grid — block commented out below, pointer left in text <<<<<<
\textit{(The total redshifted flux on a unit grid --- the full working is set out in Annexe D.)}

% \section*{The Total Flux on a Unit Grid}
% Assuming the light sources are located on a unit lattice grid in \( \mathbb{Z}^3 \), the distance to a source at \( (x, y, z) \) is:
% \[
% d = \sqrt{x^2 + y^2 + z^2}.
% \]
% 
% The flux from a single source, incorporating Hubble’s Law, becomes:
% \[
% F_{\text{obs}}(x, y, z) = \frac{L}{(x^2 + y^2 + z^2)} \cdot \left(\frac{c}{c + H_0 \sqrt{x^2 + y^2 + z^2}}\right)^2.
% \]
% 
% The total flux observed at the grid point is the sum over all lattice points, excluding the origin:
% \[
% F_{\text{total}} = \sum_{(x, y, z) \in \mathbb{Z}^3 \setminus \{(0, 0, 0)\}} \frac{L}{(x^2 + y^2 + z^2)} \cdot \left(\frac{c}{c + H_0 \sqrt{x^2 + y^2 + z^2}}\right)^2.
% \]
% 
% 
% \begin{figure}[H]
% \includegraphics[width=1.0\linewidth]{Illustrations/vign-yang2.png}
% \caption{Yang pondering the glare from an infinitude of candles}
% \end{figure}
% 
% Lee noted two competing effects in the flux sum:
% \begin{enumerate}
%     \item \textbf{Distance Diminishes Brightness:} The inverse square law ensures that more distant sources contribute less to the total flux.
%     \item \textbf{Redshift Reduces Flux:} Hubble’s Law implies that the observed flux from distant sources diminishes further due to redshift, with the term \( \left(\frac{c}{c + H_0 d}\right)^2 \) becoming smaller as \( d \) increases.
% \end{enumerate}
% 
% To analyze the behavior of the flux, Lee approximated the redshift term for large distances:
% \[
% \frac{c}{c + H_0 d} \approx \frac{1}{H_0 d},
% \]
% leading to an approximate flux contribution:
% \[
% F_{\text{obs}}(x, y, z) \sim \frac{L}{(x^2 + y^2 + z^2)^3}.
% \]
% 
% The total flux then becomes:
% \[
% F_{\text{total}} \sim \sum_{(x, y, z) \in \mathbb{Z}^3 \setminus \{(0, 0, 0)\}} \frac{L}{(x^2 + y^2 + z^2)^3}.
% \]
% 
% This sum converges for \( s > 3/2 \), suggesting that the total flux observed is finite and dominated by contributions from nearby sources.
% 
% >>>>>> end of annexed block <<<<<<
\section*{Interpreting the Redshifted Flux}
Lee was struck by the philosophical implications of his calculations:
\begin{itemize}
    \item The finite total flux reinforced the idea that distant sources contribute less not only due to their distance but also due to the accelerating expansion of the universe.
    \item The dominance of nearby sources mirrored the observer’s limited horizon, where most of the visible light originates from the local universe.
    \item The interplay between geometry (lattice structure) and physics (redshift) revealed deeper connections between number theory and cosmology.
\end{itemize}

As Lee contemplated his findings, he marveled at the simplicity and profundity of the universe’s laws. The infinite grid of light sources, governed by Hubble’s Law, was not just a mathematical abstraction but a reflection of the universe’s expansion.

\section*{Sub-Maths: Manifesto of Merriweather circle}

\lettrine{W}{illard} leaned back in his chair, his expression thoughtful but charged with quiet conviction. Across the table, Lee Yang watched him with the measured gaze of someone who respected Willard’s intellect but was wary of his penchant for challenging orthodoxy.

\begin{figure}[H]
    \centering
    \includegraphics[width=\textwidth]{Illustrations/vign-willardYang.png}
    \caption{Willard not vibing it with Yang}
    \label{fig:pirallingprimes}
\end{figure}

“Lee,” Willard began, his voice deliberate, “you’ve always struck me as someone who isn’t afraid to think against the grain. That’s why I feel you’ll appreciate this. What I’ve been doing with number theory—these ‘sub-mathematical probings,’ as I call them—is my attempt to explore mathematics not as an immutable monolith, but as a dynamic landscape of patterns waiting to be observed, modeled, and understood.”

Yang raised an eyebrow. “Sub-mathematical? That’s a curious term. Are you suggesting the opposite of meta-mathematics?”

“Precisely,” Willard replied. “Meta-mathematics, as the Austrians and their Vienna Circle positivists conceived it, saw mathematics as an abstract, axiomatic edifice—detached from experience, immune to the flux of time and matter. Carnap even had the audacity to claim that mathematics isn’t a language for expressing thoughts but merely the logical syntax of language. To them, mathematics was a sterile game, bound by rules but devoid of meaning.”

Yang nodded slowly, his gaze sharpening. “And Gödel proved them wrong.”

“Exactly,” Willard said, his tone rising slightly. “Gödel demonstrated that even within their precious axiomatic systems, there would always be truths that couldn’t be proven—truths that transcended their mechanical frameworks. And here’s the irony: Gödel, this master of abstraction, also discovered a static solution to Einstein’s field equations—a timeless, immutable picture of spacetime that defied the mutability of the physical universe.”

Yang interjected, “A picture that banished time itself. A \textit{Weltbild}—a worldview of reality as eternal, unchanging.”

“Right!” Willard exclaimed. “But where Gödel sought immutability, I see the opposite in mathematics. The objects of set theory may appear timeless, as Gödel claimed, but mathematics itself evolves—its truths, its tools, its relevance to science. This is where I diverge from the axiomatic purists. I don’t see mathematics as a temple to chip away at, stone by stone, with logical chisels. I see it as data—a sprawling, infinite dataset that demands exploration, visualization, and pattern recognition.”

Yang studied him for a moment, then asked, “But why call it sub-mathematics? Doesn’t that diminish its significance?”

Willard shook his head. “Not at all. It’s sub-mathematics because it grounds itself in the mutable, the empirical, the temporal. Where meta-mathematics floats above, detached and axiomatic, sub-mathematics delves below—into the chaotic, messy data of reality. It embraces the imperfections of computation, the contingency of models, the evolving nature of mathematical inquiry itself. It’s the complement to science, bound to time and mutability.”

Yang’s expression softened, as though he were beginning to see the contours of Willard’s vision. “You’re saying that sub-mathematics embraces the fleeting and the dynamic, while meta-mathematics chases the eternal and the unchanging.”

“Exactly,” Willard said, his voice almost reverent. “And isn’t that the ultimate irony of Gödel? In proving the limits of axiomatic systems, he paved the way for a mathematics that isn’t bound by those limits—a mathematics that can evolve, that can mirror the very flux of the universe. It’s not about rejecting immutability but complementing it with the mutable. Think of it as Planck meeting Gödel, quantum meeting continuum.”

Yang leaned back, his smile faint but genuine. “You make a compelling case, Willard. Perhaps there’s more to number theory than I gave it credit for. Maybe mathematics, like physics, needs its empiricists.”

% Willard chuckled softly. “And maybe physics, like mathematics, needs its dreamers.”

% \subsection*{Inter-Mutability in Physics and Mathematics}

Willard’s face lit up as he continued, “There’s something Roger Penrose once said about time that I think ties this all together. In physics, only objects with mass ‘experience’ time. And this experience is connected to the famous relationship between energy, mass, and time. Let me show you.”
He turned over his knapkin:
\[
E = hf= \frac{h}{T},
\]
where \(E\) is energy, \(h\) is Planck’s constant, and \(f\) is the frequency and using the relationship \(f = \frac{1}{T}\), where \(T\) is the period of oscillation. Now, equating two expressions for energy through Einstein’s mass-energy equivalence, \(E = mc^2\), and solving for \(T\) we find:
\[
mc^2 = \frac{h}{T}, \quad 
T = \frac{h}{mc^2}.
\]

“This equation,” Willard explained, “shows that the ‘experience’ of time—its duration—is tied directly to mass. Only objects with mass are subject to time’s mutability. Mathematics, in its immutable forms, doesn’t experience time. But what if we explore the mutable side of mathematics? That’s what I mean by sub-mathematics.”

Yang stared at the equation, nodding slowly. “So you’re saying that sub-mathematics brings mutability to an otherwise static discipline. A kind of temporal mathematics.”

“Exactly,” Willard said. “And in doing so, it gives mathematics a way to relate to the physical universe—not just as a tool, but as a dynamic partner in understanding reality.”

% \subsection*{Sub-Mathematics: A Grounding Manifesto}

Yang leaned back, the faint flicker of amusement crossing his face. "You know, Willard, if we're coining terms for this... perspective of yours, why not dabble with something a bit more... poetic? We've got \textit{meta-}, which means 'beyond' or 'about,' so maybe the antonym could be something like \textbf{intra-}, to signify 'within' mathematics."

Willard smiled faintly, shaking his head. "Intra- is tempting, but it feels... enclosed, doesn't it? This isn't about being constrained. It's about diving beneath the surface to what underlies mathematics. Something deeper."

Lee tilted his head. "Fair. Then what about \textbf{infra-}? It has that sense of being beneath, like in 'infrastructure.' A hidden, supporting layer."

Willard paused thoughtfully but shook his head again. "Close, but 'infra-' feels... functional. Like scaffolding or a wavelength band. We're not just dealing with hidden mechanics here; we're talking about a grounding philosophy—something that underscores mathematics itself."

Lee tapped his pen against the edge of his notebook. "How about \textbf{hypo-}? It has that same sense of 'beneath' but a more conceptual feel, like 'hypothesis' or 'hypothetical.'"

Willard gave a soft laugh. "Hypo- sounds too tentative, like we're theorizing. This isn't about speculation. It's about \textit{what mathematics stands on}, not some conjectural layer."

Finally, Willard leaned forward, his tone turning resolute. "\textbf{Sub-} is the word. It's not just beneath—it grounds mathematics. It’s the opposite of meta- in the clearest sense. Where meta-mathematics abstracts away from mathematics to examine its structure, \textbf{sub-mathematics} delves into its foundations. Think of mappings in mathematics—\textbf{injective} and \textbf{surjective}. Injection maps \textit{into} a structure, while surjection -as the French\footnote{French \textit{sur}=on of English.} would have us know- maps \textit{onto} it. Sub-mathematics is like an 'injection' into the grounding truths beneath mathematics itself."

Lee nodded slowly. "I can see that. \textbf{Sub-} has a clarity the others lack. It’s not about abstraction, it’s about anchoring. A kind of... foundation-laying for numbers."

% Willard smiled, pleased. "Exactly. Let’s not build castles in the air, Lee. Let’s anchor them in the bedrock of the earth."


\chapter*{Nambu's Prime Shadows of Composites}

\lettrine{Y}{oshiro} Nambu, a mathematician with a keen climber of snowy Himalayan peaks and no better known for interest in the interplay between primes and composites, often found himself contemplating the nature of accumulation. 

\begin{figure}[H]
\includegraphics[width=1.0\linewidth]{Illustrations/vignette_nambu.png}
\caption{Nambu contemplating his accumulated flux}
\end{figure}


For him, composites were not merely the products of primes but dynamic entities shaped by the flux of their prime factors. Viewing each composite as an "accumulated flux" of contributions from its nearest prime factors, Yoshiro sought to explore how these contributions diminished as composites grew larger. His ideas began to echo those of Willard Jones, who had previously explored the flux of light reaching an observer from lattice-based sources.

% >>>>>> ANNEX D (cut-sheet v2): composites as accumulated flux (merge with the Jinyoung material) — block commented out below, pointer left in text <<<<<<
\textit{(Composites as accumulated flux (merge with the Jinyoung material) --- the full working is set out in Annexe D.)}

% \subsection*{Composites as Accumulated Flux}
% Yoshiro envisioned a composite number $C$ as the result of flux contributions from its prime factors $p_1, p_2, \ldots, p_k$. Each prime $p_i$ contributed a share to $C$, inversely proportional to the size of $p_i$ and modulated by the thinning density of primes as $C$ increased and was modeled as:
% \[
% F(p_i, C) = \frac{1}{p_i \cdot \ln(C)},
% \]
% where:
% \begin{itemize}
%     \item $F(p_i, C)$ represents the flux contribution of $p_i$ to $C$,
%     \item $p_i$ is a prime factor of $C$,
%     \item $\ln(C)$ reflects the density of primes thinning as $C$ increases.
% \end{itemize}
% 
% The total flux $F(C)$ for a composite $C$ was then given by:
% \[
% F(C) = \sum_{p_i \mid C} \frac{1}{p_i \cdot \ln(C)}.
% \]
% 
% \section*{Connecting to Willard Jones' Lattice Flux}
% Willard’s exploration of light flux from lattice sources had emphasized the diminishing contribution of distant sources due to the inverse square law. Yoshiro extended this idea to the prime shadows, viewing primes as "sources of arithmetic light" whose flux contributions diminished both due to their increasing magnitude and their thinning density.
% 
% Instead of spatial attenuation, Yoshiro's primes experienced "numerical attenuation" governed by the inverse relationship with $\ln(C)$, paralleling the cosmological dimming of Willard’s distant light sources under Hubble’s law.
% 
% \section*{Prime Multipliers and Composite Growth}
% Yoshiro noted that composites, while products of primes, grew at a slower rate relative to the accumulation of all primes. Unlike the factorial $n!$, which grows explosively and leaves no integers untouched in its interval, the composites had gaps determined by their prime structure. He had in mind the classic proof of the infinitude of primes, which relies on factorials:
% \[
% n! + 1 \quad \text{is either prime or divisible by a prime not in } \{2, 3, \ldots, n\}.
% \]
% In contrast, the composite numbers in Yoshiro’s framework were defined by their "prime multipliers," or the set of primes contributing to their formation. If $C$ is a composite formed by primes $p_1, p_2, \ldots, p_k$, the prime multiplier factorial $P(C)$ is their \textbf{radical}:
% \[
% P(C) = p_1 \cdot p_2 \cdot \ldots \cdot p_k.
% \]
% For example:
% \begin{itemize}
%     \item For $C = 60 = 2^2 \cdot 3 \cdot 5$, $P(C) = 2 \cdot 3 \cdot 5 = 30$ (ignoring multiplicity).
% \end{itemize}
% 
% \subsection*{The Diminishing Flux of Primes}
% Yoshiro was fascinated by the interplay of diminishing contributions from larger primes as composites grew larger. The flux contributions $F(p_i, C)$ weakened not only because of the inverse proportionality to $p_i$ but also because the density of primes thinned as $C$ increased.
% 
% To quantify this, he explored the Prime Number Theorem:
% \[
% \pi(x) \sim \frac{x}{\ln(x)},
% \]
% where $\pi(x)$ counts the number of primes less than or equal to $x$. For a composite $C$, the density of contributing primes was proportional to:
% \[
% \frac{\pi(C)}{C} \sim \frac{1}{\ln(C)}.
% \]
% Thus, the thinning of primes introduced an additional decay factor to the total flux:
% \[
% F(C) \sim \frac{1}{\ln(C)} \sum_{p_i \mid C} \frac{1}{p_i}.
% \]
% Yoshiro referred to the flux contributions from prime factors as "prime shadows" cast on the composite number. These shadows dimmed as:
% \begin{enumerate}
%     \item The composite grew larger, increasing \(\ln(C)\),
%     \item The contributing primes became sparser, reducing \(\sum_{p_i \mid C} \frac{1}{p_i}\).
% \end{enumerate}
% 
% The metaphor of prime shadows highlighted the hidden influence of primes on composites, shaping their structure while diminishing their individual presence in the total flux. 
% 
% Yoshiro often reflected on the philosophical connections between her prime shadows and Willard’s light flux:
% \begin{itemize}
%     \item Both frameworks explored how individual sources—whether primes or lattice-based light emitters—contributed to a cumulative entity.
%     \item The diminishing flux, governed by $\ln(C)$ or the inverse square law, revealed the hidden structures underlying apparent randomness.
%     \item The composite numbers, much like Willard’s observer’s lattice, served as a canvas for the interplay of diminishing contributions.
% \end{itemize}
% 
% As Yoshiro closed his notebook, he thought about the beauty of numbers: how primes, the building blocks of arithmetic, shaped composites through their diminishing flux. For him, the study of composites was not just about numbers but about uncovering the hidden patterns and connections that defined the mathematical universe.
% 
% >>>>>> end of annexed block <<<<<<
\section*{Nambu \& Yang on Multipole Expansions}

Koshiro Nambu adjusted his glasses, his gaze fixed on the notebook in front of him. The handwritten equations—his own—seemed out of place in the polished, orderly realm of pure mathematics. Across from him, Lee Yang sat with a reserved expression, his arms folded, as if waiting for Nambu to justify this foray into unfamiliar territory.

Nambu cleared his throat. “Yang, I’ve been thinking about multipole expansions. You’re familiar with the concept, I assume?”

Yang tilted his head slightly, a hint of skepticism in his eyes. “In physics, yes. The decomposition of fields into monopole, dipole, quadrupole components—descriptions of charge or mass distributions. But I wasn’t aware it had anything to do with number theory.”

“It doesn’t. At least, not yet,” Nambu admitted. “But hear me out. In physics, the multipole expansion isn’t just a descriptive tool. It’s a way of organizing complexity. The monopole gives you the gross total—say, the net charge. The dipole adds directional information. The quadrupole goes further, describing asymmetries. And so on.”

Yang raised an eyebrow. “Go on.”

Nambu leaned forward, his voice steady but deliberate. “You deal with primes, don’t you? Isolated, singular entities. But they aren’t always isolated, are they? You have twin primes—\( (p, p+2) \)—and cousin primes—\( (p, p+4) \). There’s structure there, patterns that seem to form groupings. What if those groupings are like the moments in a multipole expansion?”

Yang’s posture softened slightly, though his skepticism lingered. “You’re suggesting primes can be categorized like charges in a field?”

\begin{figure}[H]
\centering
\includegraphics[width=0.8\linewidth]{Illustrations/vign-nambuStare.png}
\caption{Nambu picturing the Moments of the CMBR.}
\end{figure}

“Something like that,” Nambu said. “Take the twin primes. Two entities, closely linked, like a dipole moment. Cousin primes, separated by \( 4 \), could correspond to a quadrupole. Sexy primes—\( (p, p+6) \)—might be the octupole. The further the separation, the finer the detail.”

Yang frowned, his mind turning over the idea. “And what’s the monopole in this analogy? A single prime?”

“Perhaps,” Nambu said, pausing. “Or perhaps the overall density of primes. The distribution itself could be seen as the monopole, with deviations—these groupings of twins, cousins, quadruples—acting as higher-order moments.”
\smallskip

Yang tapped a pen against the table, his skepticism giving way to intrigue. “It’s... an interesting analogy. But in mathematics, we’re less concerned with descriptive hierarchies. We want precise relationships, proofs, something concrete to work with.”
\smallskip

“I’m not claiming it’s rigorous,” Nambu said, his tone measured. “But the idea has merit. Consider the fields generated by charge distributions. Multipole expansions describe those fields in detail, breaking them into components. Could primes generate their own kind of ‘field’? Something that dictates their groupings?”
\smallskip

Yang’s pen stilled. “You mean like a potential field? A mathematical construct that influences how primes behave?”
\smallskip

Nambu nodded. “Exactly. And if such a field exists, perhaps it’s tied to the patterns you see—the twin primes, the quadruples. Even the higher groupings, like the prime quadruples \( (p, p+2, p+6, p+8) \). Each level could be a kind of moment, describing the geometry of the number line.”
\smallskip

Yang glanced down at his own notes, a reluctant curiosity taking hold. “It’s a stretch, but not an uninteresting one. There’s something satisfying about the idea of primes forming higher-order structures. But I’m not sure I see the connection to the zeta function or number theory as a whole.”
\smallskip

“I’m not pretending it’s a direct connection,” Nambu said, his voice firm but not defensive. “But physics has a way of wandering into mathematics. Look at quantum field theory—entirely dependent on the language of pure math. This multipole analogy might seem naïve to you, but I think it could open doors.”
\smallskip

Yang allowed himself a small smile. “So this is you wandering into pure mathematics? It suits you, Nambu. But if you’re serious about this, you’ll need more than an analogy. What would a ‘multipole expansion’ of primes even look like?”
\smallskip

Nambu hesitated, aware of the thin ice he was treading on. “That’s what I’d like to explore. It’s just an idea for now. But you’ve been studying primes for decades. Surely you’ve seen patterns that resist simple classification. Maybe this framework could help.”
\smallskip

Yang leaned back, his skepticism not entirely gone but tempered by curiosity. “Maybe. But I’d need more than a vague analogy before I’d take it seriously. Still, it’s an idea worth mulling over. Prime numbers as moments... who would have thought?”
\smallskip

Nambu smiled faintly, a quiet satisfaction in his expression. “Sometimes it takes an outsider to see things differently. Physics has been moving closer to pure math for years now. This might be one more step.”
\smallskip

Yang glanced at the clock, then back at Nambu. “It’s an interesting thought experiment, I’ll give you that. But don’t expect me to do the heavy lifting. If you’re going to wander into my territory, Nambu, you’d better be prepared to stay a while.”
\smallskip

Nambu muttered to himself indiscernibly, sound low and measured. “Fair enough, Yang. But let’s me just show you want we have on lonely primes to be sure we are talking the same language.”
\smallskip

Yang leaned forward, clearing his throat a little conscious he was stating the obvious but without too much care he was infantalising Nambu who had never revealed any signs of being patronised  . 
\smallskip

``The study of prime distributions often emphasizes coupled prime pairs, such as \textbf{twin primes}—pairs of the form 
\((p, p+2)\). While twin primes have garnered significant attention, the larger question is how such pairings evolve as gaps increase.''
\smallskip

``The natural conjecture is that as \(p\) grows larger, twin primes become less frequent because the likelihood of both \(p\) and \(p+2\) being prime diminishes. Approximately, the probability of a randomly chosen odd number \(p\) being prime is about \(\frac{1}{\ln(p)}\), so the joint probability of \((p, p+2)\) being prime is roughly:
\[
\frac{1}{\ln(p) \cdot \ln(p+2)}.
\]''
Nambu nodded earnestly and poltitely. 

\smallskip

``However,'' Yang continued, ``we can extend this definition to generalize prime pairs as \((p, p+2n)\) for any integer \(n\). As the gap \(2n\) increases, the probability that at least one number in \([p, p+2n]\) is composite also rises. These pairs are often categorized as follows:''
\begin{itemize}
    \item \textbf{Cousin Primes:} Pairs \((p, p+4)\), with \(n=2\).
    \item \textbf{Sexy Primes:} Pairs \((p, p+6)\), with \(n=3\).
    \item \textbf{Generalized Pairs:} Any \((p, p+2n)\) for larger \(n\).
\end{itemize}
\smallskip

``We can just as well define coupled primes of the form ${p,p+2n}$ for any integer n. For larger gaps, \(n > 1\) the term "twin primes" is replaced by more general descriptions, such as "cousin primes" for \(n=2\) and "sexy primes" for \(n=3\), or simply "primes with a gap of \(2n\)" for larger \(n\)."
\smallskip

"We can reasonably expect the frequency of these generalized prime pairs to decrease more rapidly than that of twin primes as \(n\) increases. This is because the larger the gap \(2n\), the higher the probability that at least one number within the interval \([p, p+2n]\) is composite, a likelihood that escalates with \(n\).''

\begin{figure}[H]
    \centering
    \includegraphics[width=0.9\textwidth]{Illustrations/twincousin.png}
    \caption{Ulam Spiral highlighting cousin Prime-pairs up to 940 (\(n=2\), gap=4).}
    \label{fig:cousinprimes}
\end{figure}

''Yes,'' Nambu interjected, ''The code, \href{https://colab.research.google.com/drive/1qyDEbWt6gnFmL75O_J7vwe2K3pLWXhcg?usp=sharing}{TwinPrimeUlamSpiral-plot-anyTwins.ipynb}, delivers a plot of paired cousin primes, where the gap between primes is \(4\).

\begin{figure}[h!]
    \centering
    \includegraphics[width=0.6\textwidth]{Illustrations/twinprime10000.png}
    \caption{Ulam Spiral highlighting twin Prime-pairs up to 10000 (\(n=1\), gap=2).}
    \label{fig:twinprimes}
\end{figure}
\noindent
Here the focus is on twin primes (\(n=1, \text{gap}=2\)), with green rectangles marking the pairs and coloring represents primes modulo 4:\textbf{Red dots:} Primes of  form \(4n+1\); {Blue dots:} Primes of the form \(4n+3\).

\begin{figure}[h!]
    \centering
    \includegraphics[width=0.6\textwidth]{Illustrations/singltonprimeToorder8100000.png}
    \caption{Ulam Spiral Highlighting Lonely Primes up to 100000.}
    \label{fig:singletonprimes}
\end{figure}

The final plot identifies \textbf{lonely primes}, which are not part of any pair with gaps \(2, 4, 6, 8, 10, 12, 14, 16\). These primes, shown with black rectangles, highlight the \textit{hollowing out} phenomenon in the spiral—composite numbers increasingly fill gaps as numbers grow.

Yang studied the plots silently, nodding. \textit{``So, as we increase \(n\), pairs become sparse, and singletons dominate?''}

\textit{``Exactly,''} Nambu replied. ``The patterns suggest that as numbers grow, composites flood the gaps, isolating primes into rarefied regions.''

\begin{figure}[ht]
    \centering
    \includegraphics[width=0.8\textwidth]{Illustrations/vign-nambuSpiral.jpg}
    \caption{Nambu spiraling into a depression.}
    \label{fig:pirallingprimes}
\end{figure}


Nambu pointed to the singleton plot. \textit{``This hollowing out—doesn’t it reflect some underlying structure, like a multi-moment expansion? Charge distributions fragmenting into monopoles, dipoles...''}

Yang allowed a faint smile. ``Perhaps. But primes, like physics, often resist neat analogies. That’s the beauty of it, isn’t it?''

% Yang leaned back. \textit{``Beauty, or chaos? I suppose it depends on your perspective.''}
%%%%%%%%%%%%%%%%%

\section*{Prime Spiral Polar Rays: A Discussion}

Nambu returned to his office, the weight of the day’s reflections pressing down on him. The discouraging patterns of lonely primes, fading away as higher-order \(n\) pairs filled the gaps, left him in a somber mood. He slumped into his chair, his mind drifting between the pages of his latest calculations and the thought that perhaps the primes were leading him into a wilderness with no end.

The phone on his desk buzzed, breaking his reverie. He hesitated, then picked it up. \textit{``Nambu here.''}

\textit{``Koshiro, it’s Willard Jones. I hope I’m not catching you at a bad time,''} Willard said, his tone light but with a hint of concern.

\textit{``Willard,''} Nambu replied, trying to inject some energy into his voice. \textit{``What can I do for you?''}

\textit{``Actually, it’s more about what I can do for you,''} Willard said. \textit{``I heard you’ve been grappling with the lonely primes.''}

Nambu leaned back, rubbing his temples. \textit{``I have. It’s as if they vanish into the ether as the higher-order pairs take over. I’m starting to wonder if this whole line of inquiry is a dead end.''}

Willard’s voice softened. \textit{``Nambu, don’t lose hope just yet. Let me draw your attention to a different perspective. Think of the lonely primes not as failures of the system, but as its monopoles—singular entities in a universe of binaries and higher-order structures. There’s beauty in their rarity.''}

Willard nodded. \textit{``Precisely. Consider stars in the cosmos. Binary systems dominate, yet it’s the solitary stars—the singular, wandering luminaries—that capture our imagination. Lonely primes are the monopoles of the numerical universe, rarer, but no less significant.''}

\textit{``What do you mean?''} Nambu asked, intrigued despite himself.

\textit{``Let me show you something,''} Willard said. \textit{``Prime numbers, except for 2 and 3, conform to the residue classes \(6n \pm 1\) to avoid divisibility by 2 or 3. On a polar plot, these primes form spirals, clustering at angles corresponding to these forms. The lonely primes—those not part of any pair—are like monopoles in the cosmic fabric of numbers.''}

Nambu studied the plot Willard had sent over email.
\begin{figure}[ht]
    \centering
    \begin{subfigure}{\textwidth}
        \centering
        \includegraphics[width=\linewidth]{Illustrations/logSpiral6n+11000000dark.png}
        \caption{Polar spirals delineating prime types.}
        \label{fig:z6n+1Spiral}
    \end{subfigure}
    \caption{Prime residue patterns in polar coordinates.}
    \label{fig:both}
\end{figure}

Willard continued, \textit{``Similarly, primes in residue classes mod 4 (\(4n+1\) and \(4n+3\)) form distinct patterns across rays. These spirals aren’t just random—they’re deeply tied to the structure of numbers and Euler’s totient function.''}

\textit{``Euler’s totient function?''} Nambu asked, now fully engaged.

\textit{``Exactly,''} Willard said. \textit{``Take a number like 710. Its prime factorization is \(710 = 2 \cdot 5 \cdot 71\). Euler’s totient function, \(\phi(710)\), counts 280 integers coprime to 710. These coprime integers correspond to the missing teeth in the radial distribution of primes on the spiral.''}

\section*{Twin Primes and Higher-Order Pairs}

\textit{``And pairs?''} Nambu asked, leaning forward.

\begin{figure}[h!]
    \centering
    \includegraphics[width=0.9\textwidth]{Illustrations/twinPrimes100001.png}
    \caption{Spiral of twin primes up to \(10^6\).}
    \label{twinSpiral:fig}
\end{figure}


''Ah, pairs are where the structure really gets interesting,'' Willard said. ``Twin primes, like \((p, p+2)\), form the most basic pairs. But as you know, we can generalize to cousins \((p, p+4)\) or even higher gaps. These pairs appear less frequently as \(n\) increases, but their patterns remain. Here in polar coordinate scatterplots where now each point, $(n,n)$ instead represents an integer $n$ up to $10^6$, and the radial coordinate is proportional to $\log_{10}(n)$. Twin primes appear in 9 pairs across each quadrant.

\section*{Lonely Primes: The Monopoles of Number Theory}

Willard paused, then gestured toward another plot. ``These are the lonely primes, Nambu-pairless. They are hollowed out from the spiral, as composites fill the void.''

\begin{figure}[H]\centering
\includegraphics[width=\textwidth]{Illustrations/lonelyPrimes100000_16.png}
\caption{Lonely primes amongst first 1000000 up to order 16.}
\label{twinSpiral:fig}
\end{figure}

Nambu leaned back, his expression lighter. \textit{``You’ve given me something to think about, Willard. The primes—lonely or not—are part of a much larger story.''}

Willard nodded. \textit{``And as with monopoles, their rarity doesn’t diminish their importance. If anything, it enhances it.''}

% Nambu allowed himself a rare smile. \textit{``Perhaps there’s more to the spiral than I thought.''}

%%%%%%%%%%%%
\section*{Piecewise Analysis of Lonely Primes}

Willard leaned in closer to the screen as Nambu shared his findings.

Nambu adjusted his glasses and began, his tone deliberate. "Willard, what you’re looking at here are the lonely primes—primes that are not part of any pairs within a specified gap."

% \subsection*{Table of First 100 Lonely Primes (up to $n=2$)}
\begin{longtable}{|c|c|c|c|c|}
\hline
\textbf{Prime 1} & \textbf{Prime 2} & \textbf{Prime 3} & \textbf{Prime 4} & \textbf{Prime 5} \\ \hline
53  & 89  & 157 & 173 & 211 \\ 
251 & 257 & 263 & 293 & 331 \\ 
337 & 359 & 367 & 373 & 389 \\ 
409 & 449 & 479 & 509 & 541 \\ 
547 & 557 & 563 & 577 & 587 \\ 
593 & 607 & 631 & 653 & 683 \\ 
691 & 701 & 709 & 719 & 727 \\ 
733 & 751 & 787 & 797 & 839 \\ 
919 & 929 & 947 & 953 & 977 \\ 
983 & 991 & 997 &      &     \\ \hline
\caption{Table of First 100 Lonely Primes (up to $n=2$)}
\end{longtable}

Nambu gestured to the table. "This first table shows lonely primes who have no cousins. As the gap increases, the criteria for pairing expand, and the pool of lonely primes shrinks." Nambu pointed to the second table. "Here, we extend the maximum pairing gap to \(2n = 16\)."

\begin{table}[h]
\centering
\begin{tabular}{|c|c|c|}
\hline
\textbf{Prime 1} & \textbf{Prime 2} & \textbf{Prime 3} \\ 
\hline
2179            & 2503            & 3967            \\ 
4177            & 7369            & 7393            \\ 
\hline
\end{tabular}
\caption{Table of Lonely Primes for the First 10,000 Numbers}
\label{tab:lonely_primes}
\end{table}


Willard raised an eyebrow. "The primes here seem more isolated. Is there a threshold beyond which no primes remain lonely?"

Nambu nodded. "That’s the part I am looking into. I am just running some python code now over the first 10 million numbers to see which prmes do not have a pair within 64 places. Naturaly as the pairing gap \(2n\) grows, lonely primes become increasingly rare."

Willard leaned back. "So these tables are snapshots, but the full story lies in the trend?"

"Exactly," Nambu replied. "And that’s what I am plotting at the moment."

Willard’s voice came through the call, cutting through Nambu’s concentration. "Koshiro, what are you seeing there? Looks like something’s eating at you."


\begin{figure}[H]
    \centering
    \includegraphics[width=\textwidth]{Illustrations/piecewiseLonelyPrime.png}
    \caption{Piecewise Model Fits for \(m\) vs \(2n\): The plot shows the number of lonely primes \(m\) as a function of the maximum pairing gap \(2n\). The data (blue points) is fit with an exponential decay model (orange dashed line) for \(2n \leq 34\), transitioning to a logarithmic decay model (green dashed line) for \(2n > 34\). The red dashed line represents the piecewise model combining both fits.}
    \label{fig:piecewiseLonelyPrime}
\end{figure}

The piecewise model is defined as:
\[
m =
\begin{cases} 
758617 \cdot \exp(-0.15 \cdot 2n), & \text{for } 2n \leq 34, \\
-145000 \cdot \log(2n) + 537290, & \text{for } 2n > 34.
\end{cases}
\]

Nambu sat silently for a moment, staring at the diminishing trend of lonely primes on his screen. he then shared his screen, gesturing to the plot in frustration. "This is the analysis of lonely primes as a function of \(2n\), the maximum pairing gap.

% Initially, the decay is sharp, exponential but beyond that, it transitions to a logarithmic decay:

Willard leaned in, studying the plot. "A piecewise fit then. Makes sense. The early exponential decay reflects the rapid thinning of lonely primes as gaps widen. But as the gaps grow larger, the system stretches into the logarithmic regime—a slower but more persistent decline."

\bigskip

Nambu nodded but looked unconvinced. "Yes, but what drives this transition? Why do lonely primes follow two distinct behaviors?"
\bigskip

Willard chuckled. "Maybe there’s beauty in the change itself. The exponential decay mirrors the collapse of simpler structures borne of the greater density of primes at the beginning, while the logarithmic tail hints at a hidden coherence—like the primes themselves resisting oblivion."
\bigskip

Nambu raised an eyebrow. "Maybe we don't need to understand the mathematics behind the transition point at \(2n = 34\). What shifts at that scale?"

% Willard gestured to the screen. "Let’s look at the plot and see if it offers any clues."

Willard leaned back. "Early on, the lonely primes vanish rapidly, but as gaps grow, the system stabilizes into a slower decline. This asymptotic behaviour is what we are reall interested- when the primes are settling into an attenuating equilibrium."
\bigskip

Nambu tapped his pen thoughtfully. "Perhaps. Or perhaps the transition reflects a shift in the arithmetic structures governing these primes. Residue classes, modular symmetries."
\bigskip

Willard grinned. "Well, Koshiro, for now, take solace in the fact that these primes, lonely as they are, refuse to fade away entirely. There’s beauty in their persistence."
\bigskip

Nambu allowed a faint smile. "Persistence indeed. Let’s see where that leads us."


% \section*{postscript: Polar cycle}
% \href{https://www.youtube.com/c/3blue1brown}{3Blue1Brown} in a  \href{https://www.youtube.com/watch?v=EK32jo7i5LQ}{youtube video} explores the prime spoke behavior that conjures from the inherent spiral plotting of $(n,n)$ on a non-complex polar plane as a visual consequence of the relationship between \textbf{Euler's totient function}, $\phi$ and continued fraction approximations of $\pi$. As we are mapping primes onto polar coordinates, their angular position is proportional to their magnitude. Since $3/1$ is a crude approximation of $\pi$, we find that $6n$ positions map to approximately $2\pi$, meaning that after six "turns", the points have traversed just under a full circle.  Improved approximations can be obtained from the continued fraction representation of $\pi$:
% \begin{align*}
% \pi = [3; 7, 15, 1, 292, \ldots] & \approx 3 + \cfrac{1}{7 + \cfrac{1}{15 + \cfrac{1}{1 + \cfrac{1}{292 + \cdots}}}}\\
% % &\approx 3 + \frac{1}{7} = \frac{22}{7}\ \ ;\ \ 3 + \frac{1}{7 + \frac{1}{15}} = \frac{333}{106}\ \  ;
% &\approx 3 + \frac{1}{7 + \frac{1}{15 + \frac{1}{1}}} = \frac{355}{113}
% \end{align*}
% \noindent
% A better approximation, $\pi\approx 22/7$, implies that $44n$ positions correspond to nearly $2\pi$, creating seven almost complete turns and leads to the primes $4n \pm 1$ delineating the circular structure. $\frac{355}{113}$ as the best rational approximation of $\pi$ with a denominator of three digits results in 113 turns delivering an almost completely divergent ray structure with barely a hint of spiralling. 710 as twice the numerator of the fraction $\frac{355}{113}$ has prime factorization $710 = 2 \times 5 \times 71$, and is divisible by 5 which we observe in the polar plot as the primes come in batches of 4. 

% Euler's {\em totient function}, $\phi(n)$ counts the integers up to $n$ that are coprime (relatively prime) to $n$. We have $\phi(710)= 710\cdot\frac{1}{2} \cdot \frac{4}{5} \cdot \frac{70}{71}=280$ filtering out multiples of 5 giving the "missing teeth" in the radial distribution. The terms $\frac{1}{2}$, $\frac{4}{5}$, and $\frac{70}{71}$ act as filters to exclude non-coprimes.

\section*{A Bayesian Perspective on Loneliness}

\lettrine{E}{dward} Willard fumbled with his tray as he spotted Jovannah in the corner of the cafeteria. Her poise and ease always unnerved him, yet after his conversation with Nambu, he felt compelled to seek her insights. Jovannah, already halfway through a cup of tea, noticed his hesitation and gestured with the slightest of a come hither cupped hand for him to join her." 

\begin{figure}[h!]
    \centering
    \includegraphics[width=0.9\textwidth]{Illustrations/vign-willardJovannah.png}
    \caption{Willard a little nervy around Jovannah}
    \label{fig:pirallingprimes}
\end{figure}

Willard hesitated for a moment- more in feigned impudence as he knew he was being predated upon, before making his way to her table. Jovannah, enjoying his awkwardness, greeted him with a dutiful nod. Willard gathering some measure of control, placed a set of printed plots on the table and began, "Jovannah, I wanted to get your thoughts on something. I’ve been looking at this polar plot of twin primes of late."
\bigskip

Jovannah leaned forward, examining the plot. Before Willard could launch into his explanation, she interrupted, "Ah, yes! This reminds me of a \href{https://www.youtube.com/c/3blue1brown}{3Blue1Brown video} I saw recently." She tapped the table lightly, as if recalling the details. "\href{https://www.youtube.com/watch?v=EK32jo7i5LQ}{They explore} the prime spoke behavior using polar coordinates. Have you seen it?"
\bigskip

Willard shook his head, slightly deflated. Jovannah continued, “It’s fascinating. When you map primes onto a non-complex polar plane, their angular position becomes proportional to their magnitude. This reveals patterns—spokes, really—that arise as a visual consequence of the relationship between \textbf{Euler's totient function}, \( \phi(n) \), and continued fraction approximations of \( \pi \). For example, consider the crude approximation of \( \pi = 3/1 \): it suggests that \( 6n \) positions map to roughly \( 2\pi \), meaning that after six turns, the points almost complete a circle.”
\bigskip

Willard scratched his head nervily, trying not to look unsure. “Six turns- oh I see?”
\bigskip

Jovannah reading between the scratch and the question mark spared ambiguity and ploughed on. “So as we know, better approximations of \( \pi \) come from its continued fraction representation:
\[
\pi = [3; 7, 15, 1, 292, \ldots] \approx 3 + \cfrac{1}{7 + \cfrac{1}{15 + \cfrac{1}{1 + \cfrac{1}{292 + \cdots}}}}.
\]
For instance:
\[
\pi \approx 3 + \frac{1}{7 + \frac{1}{15 + \frac{1}{1}}} = \frac{355}{113},
\]
is the best rational approximation of \( \pi \) with a denominator of three digits. We observe remarkable behavior, with 113 turns, you get almost completely divergent rays. This refinement explains why 710—twice the numerator of \( \frac{355}{113} \)—plays a role. Its prime factorization, \( 710 = 2 \times 5 \times 71 \), ties directly to the prime batches observed in the spiral.”
\bigskip

Willard leaned back, furrowing his brow. “So, the radial positions of primes aren’t just random—they align based on modular behavior?”
\bigskip

“Exactly!” Jovannah said with a smile. “Euler’s \emph{totient function}, \( \phi(n) \), helps clarify this. It counts the integers up to \( n \) that are coprime to \( n \). For example, \( \phi(710) = 710 \cdot \frac{1}{2} \cdot \frac{4}{5} \cdot \frac{70}{71} = 280 \). These terms filter out multiples of small primes, giving the ‘missing teeth’ in the radial distribution.”
\bigskip

Willard, still digesting this, glanced at his plot again. “It is probably worth looking at the number of radial for other prime couplings beyond the twins to see their connection to this modular framework.”

% \subsection*{Bayesian Analysis of Lonely Primes}

Willard cleared his throat as he sat. “But Nambu’s work on lonely primes, it is a little vacuous perhaps?”
\bigskip

Jovannah raised an eyebrow. “Well, Willard, if we’re serious about understanding lonely primes, we should really be framing it probabilistically. Let’s start with the basics.”

% \subsection{Prior: Probability that a Number is Prime}
\bigskip

Jovannah continued, her voice steady and confident. “This is where Bayesian reasoning comes into play. To predict the likelihood of a number being a lonely prime, we combine a prior and a likelihood. The prior probability follows the Prime Number Theorem:
\[
P(\text{prime}) \approx \frac{1}{\ln p}.
\]
This captures how primes become sparser as \( p \) increases.”
\bigskip

Willard nodded, still the tone of her <<Willard>> resonating. “Right. So that’s the baseline. What about the probability that a prime is lonely?”

% >>>>>> ANNEX C (cut-sheet v2): the likelihood computation for lonely primes — block commented out below, pointer left in text <<<<<<
\textit{(The likelihood computation for lonely primes --- the full working is set out in Annexe C.)}

% \subsection{Likelihood: Probability \( p \) has no neighboring prime in its orbit}
% 
% Jovannah smiled. “That’s where it gets interesting. To be lonely, a prime \( p \) must not have any prime neighbors within a fixed range. For example, the probability that \( p+2 \) is \emph{not} prime is given by:
% \[
% P(\neg \text{twin} \mid p) = 1 - \frac{1}{\ln p \cdot \ln(p+2)}.
% \]
% Similarly, for cousin primes (\( p+4 \)) and sexy primes (\( p+6 \)):
% \[
% P(\neg \text{cousin} \mid p) = 1 - \frac{1}{\ln p \cdot \ln(p+4)},
% \]
% \[
% P(\neg \text{sexy} \mid p) = 1 - \frac{1}{\ln p \cdot \ln(p+6)}.
% \]”
% Willard interjected, “So to enforce loneliness, we extend this reasoning to all potential neighbors within a local range, say up to 64?”
% \bigskip
% 
% Jovannah nodded. “Exactly. The probability that \( p \) has no partners within 64 is:
% \[
% P(\neg \text{any partner} \mid p) = \prod_{k \in \{2,4,6,\dots,64\}} \left(1 - \frac{1}{\ln p \cdot \ln(p+k)}\right).
% \]
% This product assumes local independence and ensures all these neighbors are composite.”
% 
% :
% \[
% P(\text{lonely} \mid \text{prime}) \approx P(\neg \text{any partner} \mid p).
% \]
% Numerical evaluation of this product provides insight into the expected frequency of lonely primes.”
% 
% >>>>>> end of annexed block <<<<<<
\subsection{Posterior: The Probability That a Given Prime is Lonely}

“Now,” Jovannah continued, “using Bayes’ theorem, we can write the posterior probability that \( p \) is lonely as:
\[
P(\text{lonely} \mid \text{prime}) = \frac{P(\text{prime} \mid \text{lonely}) P(\text{lonely})}{P(\text{prime})}.
\]
Since we are focusing only on primes, we normalize this to:
\[
P(\text{lonely} \mid \text{prime}) \approx P(\neg \text{any partner} \mid p).”
\]
% This provides a concrete framework for estimating loneliness numerically.”
\bigskip

Willard, keeping up, muttered, “And this… this quantifies a prime's loneliness?”
\bigskip

“Yes,” Jovannah replied. “Bayes’ theorem ties it together
\bigskip

Willard gathering himself eager to contribute. “So, lonely primes could be analogous to \textbf{rough primes}—numbers that avoid divisibility by smaller primes. A lonely prime avoids coupling within its local range, just as a \( P \)-rough prime avoids divisibility by elements of a small prime set \( P \).”
\bigskip

Jovannah nodded. “Ah nice. And if we factor in modular biases, such as primes of the form \( 4n+1 \) and \( 4n+3 \), we could refine our model further. Do \( 4n+1 \) primes have a higher likelihood of loneliness? That’s a question worth exploring.”

Willard jotted down notes as Jovannah outlined possible extensions:
\begin{itemize}
    \item Numerically compute \( P(\text{lonely} \mid \text{prime}) \) for various ranges of \( p \).
    \item Compare against the empirical distribution of rough primes to validate the connection.
    \item Investigate modular biases, especially \( 4n+1 \) versus \( 4n+3 \), to understand how neighborhood structures influence loneliness.
\end{itemize}

As Jovannah finished her tea, she glanced at Willard. “It’s not just about the primes themselves, Willard. It’s about the structure they inhabit—the gaps, the densities, the biases. Willard, you’re onto something. Primes, whether clustered as twins or isolated as lones, reveal deeper truths. The patterns are fundamental.”
\bigskip

Willard smiled awkwardly, still awed by her self-assurance. “Thanks, Jovannah. I’ll... think about all this.”
\bigskip

Jovannah stood, her confidence unwavering. “Don’t think too hard, Willard. Sometimes, the numbers speak for themselves.”

%%%%%%%%%%%%%

% For loneliness, we must ensure \( p \) has no nearby primes. For example, the probability that \( p+2 \) is not prime is:
% \[
% P(\neg \text{twin} \mid p) = 1 - \frac{1}{\ln p \cdot \ln(p+2)}.
% \]
% Extending this to all potential partners within a range of 64, we calculate:
% \[
% P(\neg \text{any partner} \mid p) = \prod_{k \in \{2,4,6,\dots,64\}} \left(1 - \frac{1}{\ln p \cdot \ln(p+k)}\right).
% \]


\subsection*{Willard just about Takes His Leave}

Willard, still clutching his notepad tightly, glanced down at his half-empty cup of coffee, then back at Jovannah. 

\textit{“Well, this has been... illuminating,”} he said, his voice carrying an odd mixture of sincerity and discomfort. \textit{“I’ll... uh, let you get back to whatever you were doing.”}
\bigskip

Jovannah with the faintest coquettish of head tilts, the hint of a smile playing on her lips, \textit{“You mean solving the mysteries of the universe while having tea? Don’t worry, Willard, I’ll manage.”}
\bigskip

He chuckled nervously, nearly spilling his coffee as he stood. \textit{“Right. Well, thanks again for your insights. I’ll just... mull this over. Maybe revisit the spirals. Or, you know, do something practical with all this.”}
\bigskip

Jovannah gave him a knowing nod. \textit{“Good idea. And Willard, don’t forget—you’ve got good instincts.”}
\bigskip

\textit{“Sure,”} Willard replied, backing away awkwardly, nearly bumping into a chair. \textit{“Instincts. Got it. Thanks again!”}
\bigskip

As he turned to leave, Jovannah called after him. \textit{“Oh, and Willard?”}
\bigskip

\textit{“Yes?”} he responded, spinning around too quickly, nearly colliding with another table, and quietly berating himself for his clumsiness.

She smirked. \textit{“Oh, never mind.”} 

Willard sighed, shaking his head with a nervous grin, and continued towards the exit. 

Just as he was nearly out of earshot, Jovannah added, almost offhandedly, “Oh, and if you’re really keen on making sense of it, compute the Bayesian posterior, 
\(
P(\text{lonely} \mid \text{prime}.)
\)
Numerically estimate it for primes up to a range. Compare against rough primes—identify biases and gaps in prime spacing. Simulate distributions, maybe a Monte Carlo experiment, and see how often primes are left without a close neighbor.”
\bigskip

Willard hesitated mid-step, brow furrowing. He opened his mouth, as if to respond, but then thought better of it. Instead, he simply muttered, \textit{“Monte Carlo... sure... that’s exactly what I was thinking.”} 

% Jovannah shook her head, amused, as Willard disappeared down the corridor, still scribbling in his notebook.
\bigskip

Jovannah keen to have the last word, \textit{“Don’t overthink it Willard. Sometimes, clarity finds you when you’re not looking.”}
\bigskip

Willard smiled sheepishly, nodded, and made his way out of the cafeteria, .
%%%%%%%%%%%
% ============ COMMENTED OUT (cut-sheet v2): duplicate of 'Das Ein vs. Das Man' — kept version lives in 8-6jovan.tex ============
% \section*{Das Ein vs. Das Man: A Duel of Minds}
% A dimly lit private club in Vienna. The air is thick with the scent of aged wood and burning cigars. A blackboard stands between them, half-covered in an intricate derivation on spherical harmonics and Bessel functions. A bottle of whiskey sits on the table, half-empty, the amber liquid catching the light of the low-hanging lamps. 
% 
% Jovannah swirls her drink lazily, watching the way the liquid folds and unfolds upon itself. "So tell me, gentlemen—who exactly owns the truth in mathematics? The one who discovers it first, or the one who understands it best?"  
% 
% Willard straightens slightly, his grip tightening around his glass. "Truth is not owned. It is grasped. By those willing to chase it with absolute clarity." He gestures toward the board, eyes sharp. "Rigor, derivation, proof—that is the only measure of understanding."  
% 
% Ernest leans back, smirking. "Ah, but clarity is an illusion, Willard. Mathematics is as much about intuition as it is about proof. You can’t sit around proving every step when the world moves at the speed of light. Sometimes you leap. Sometimes you trust your instincts."  
% 
% "Ah, ‘intuition.’" Willard scoffs, shaking his head. "A fine excuse for sloppiness. That’s how we get half-baked physics papers masquerading as mathematics. That’s how errors creep in, how rigor is abandoned for convenience."  
% 
% "Yes, yes," Ernest waves a hand dismissively, the whiskey in his glass sloshing slightly. "And yet, whose work has applications beyond pure abstraction? Yours, trapped in a world of ideals? Or mine, where the equations breathe in reality, shaping actual phenomena?" He taps the board. "These harmonics? They aren’t just pretty functions. They dictate how sound waves move, how galaxies form. They do something, Willard."  
% 
% Willard exhales, slow, measured. "And you think that justifies sloppy thinking? That the mere presence of applicability validates any loose approximation you see fit to take?"  
% 
% "It justifies the fact that I don’t waste my time chasing ghosts." Ernest grins. "You hide in the ivory tower because you fear being wrong. I build models because I know we must be."  
% 
% Jovannah, having watched the exchange with something between amusement and calculation, finally leans forward, her eyes glinting. "So what we have here is the eternal struggle. Das Ein and Das Man."  
% 
% Willard blinks, caught slightly off guard. "That’s not—"  
% 
% She cuts him off with a grand gesture toward him. "Willard, the idealist, the one who seeks truth without compromise, without shortcuts, who would rather be alone with his principles than risk contamination by the crowd."  
% 
% She turns to Ernest, voice lowering slightly, playful but edged. "And you, dear Ernest, the pragmatist, the man of the world, bending the rules because everyone else does, because you know truth is not some shining beacon, but something negotiated, something malleable."  
% 
% Ernest chuckles, tilting his head. "Well, when you put it like that—"  
% 
% "But tell me, gentlemen." Jovannah sets her glass down, fingers idly tracing the rim. "Who is freer? The man who holds fast to his principles and watches the world move without him? Or the man who moves with the world but no longer knows where he stands?"  
% 
% Silence.  
% 
% Willard stares at the board. Ernest swirls his glass.  
% 
% Jovannah waits a beat, then lets a triumphant little smile slip through. "What’s the matter? You both seemed so certain a moment ago."  
% 
% Willard tightens his grip on the chalk. Ernest takes another sip of whiskey.  
% 
% Neither speaks.  
% 
% 
% \newpage
% ============ end commented block ============
% >>>>>> ANNEX D (cut-sheet v2): the energy-levels-of-attraction formalism — block commented out below, pointer left in text <<<<<<
\textit{(The energy-levels-of-attraction formalism --- the full working is set out in Annexe D.)}

% \section*{Jovannah raises the Energy Levels of Attraction}
% 
% The energy levels of the hydrogen atom are determined by solving the Schrödinger equation for the Coulomb potential:
% 
% \begin{equation*}
% V(r) = -\frac{e^2}{4\pi \epsilon_0 r}.
% \end{equation*}
% 
% For a single-electron atom, the time-independent Schrödinger equation in spherical coordinates is:
% 
% \begin{equation*}
% \left[-\frac{\hbar^2}{2m} \nabla^2 - \frac{e^2}{4\pi \epsilon_0 r} \right] \psi = E \psi.
% \end{equation*}
% 
% Solving this equation with appropriate boundary conditions for bound states yields the Bohr energy levels:
% 
% \begin{equation*}
% E_n = -\frac{13.6 \text{ eV}}{n^2}, \quad n = 1,2,3, \dots
% \end{equation*}
% 
% where: \( E_n \) represents the energy of the electron in the \( n \)th energy level and \( 13.6 \) eV is the Rydberg energy, the ionization energy of hydrogen.
% 
% 
% Thus, the energy levels scale as \( 1/n^2 \), meaning successive levels get closer together as \( n \to \infty \). Helium (\(^4\text{He}\)) has two protons and two electrons. Since helium has two electrons, it is not strictly hydrogenic and does not account for electron-electron interactions., and electron-electron interactions modify the energy levels. To first order it's energy levels are modified as:
% 
% \begin{equation*}
% E_n = -\frac{Z^2 \cdot 13.6 \text{ eV}}{n^2}, \quad Z=2 \text{ for helium}.
% \end{equation*}
% A more accurate model involves multipole expansion with an electric dipole moment shifting energy levels due to charge asymmetry and quadrupole moment refining the potential shape. For helium, the dipole moment of the nucleus is zero (due to symmetric proton arrangement), but the quadrupole moment of the electron cloud contributes to fine structure shifts. Jovannah considered the reciprocals of figurate numbers:
% 
% \begin{equation*}
% E_n \propto -\frac{1}{T_n}, \quad E_n \propto -\frac{1}{S_n}, \quad E_n \propto -\frac{1}{P_n}.
% \end{equation*}
% 
% As helium’s levels deviate from strict \( 1/n^2 \) hydrogenic scaling, the corrections might follow a reciprocal figurate sequence.
% 
% 
% \begin{figure}[H]
% \includegraphics[width=1.0\linewidth]{Illustrations/residuals figurative.png}
% \caption{Elemental figurative modelling}
% \end{figure}
% 
% \newpage
% >>>>>> end of annexed block <<<<<<
\section*{An Oscillating Conversation}

A quiet café. The low hum of conversation dissolves into the background, a gentle counterpoint to Willard’s restless fidgeting. Across from him, Jovannah watches, fingers idly tracing the rim of her coffee cup, her expression unreadable but undeniably amused.

"You do realize I don’t actually need you to explain quantum mechanics to me, right?" she says, breaking the silence.

Willard blinks, clearly caught mid-thought. "I—what? No, I wasn’t—"

"Weren’t you?" She tilts her head, a slow, knowing smirk forming. "Because it sure sounded like you were about to tell me all about energy levels and wavefunctions, as if I hadn’t been staring at equations my whole life."

He clears his throat, grasping for footing. "I—no, I just assumed—"

"That because I’m a musician, I only think in terms of sound waves?" She leans in slightly, fingers now toying with his abandoned pen. "That I wouldn’t know my way around an eigenfunction?"

Before he can protest, she flips open his notebook and, with a flourish, scrawls across the open page:

\begin{equation}
E_n = - \frac{h c R}{n^2}
\end{equation}

"Bohr’s formula," she says, tapping the equation. "Electron energy levels. Energy drops in proportion to the inverse square of \(n\)."

Willard exhales, running a hand through his hair, unsure whether to feel impressed or outmaneuvered. "That’s... correct."

"And what does that remind you of?" she prompts, lifting a brow.

He hesitates, staring at the equation as if the answer is on the tip of his tongue but just out of reach.

"Come on, Professor," she teases. "You’re almost there."

Then, still holding his pen, she adds below:

\begin{equation}
\sum_{n=1}^{\infty} \frac{1}{n^2} = \frac{\pi^2}{6} = \xi(2)
\end{equation}

"Bohr’s model is just Basel’s sum in disguise," she continues, tapping the paper lightly. "The universe sings, Willard. You don’t have to tell a musician about harmonics."

He studies the page, then looks at her, caught between admiration and something else he can’t quite place. "That’s... elegant."

She tilts her head, pretending to ponder. "Oh, now it’s elegant?" A soft laugh escapes her. "A minute ago, I was being given a lecture on wavefunctions. Now suddenly, I’m a genius?"

He smirks despite himself. "Well, I wouldn’t go that far."

"You wound me," she says, placing a dramatic hand over her chest. Then, voice softer now, conspiratorial, "Schrödinger just took Bohr’s intuition and turned it into standing waves. Harmonics in probability space. And you know what that means?"

Willard, still slightly wary of where she’s leading him, tries, "That... everything is a wave?"

"That everything oscillates." She lifts her coffee, takes a slow sip, eyes twinkling. "Energy, probability, ideas..." A pause, deliberate. "Even you, right now."

He straightens slightly. "Me?"

"Oh, absolutely," she says smoothly, setting her cup down with a quiet clink. "You’re practically vibrating with nervous energy." 

He exhales, shaking his head, but he can’t help the amused chuckle that escapes him.

"Must be an eigenmode of some kind," she adds, lips curling into a smirk.

And for the first time in hours, Willard is not thinking about finance, or failure, or the ruin that looms over him.  

Only the teasing glint in Jovannah’s eyes, and the way, in this moment, she is undeniably winning.

\begin{figure}[h!]
    \centering
    \includegraphics[width=0.7\textwidth]{Illustrations/vign-backWillard.png}
    \caption{Willard watching the one that got away}
    \label{fig:pirallingprimes}
\end{figure}


%%%%%%%%%%%%%%%
\
\section*{An Isoperimetric Discussion}
Willard knew he was prone to over-explaining concepts when unsettled. He was all too soon going to be performing that routine in the company of Jovannah, as she was sacheting through the door to a quiet cafe. Willard, visibly anxious, seeing her coming a little way off, not so idly pretends to doodle something akin to a spherical harmonic on a napkin while Jovannah, dressed in an off-shoulder blouse, catches his little routine with mild amusement as she flounces in beside him with not a word or a nod. A pause opens up as Jovannah lifts her eyebrows and purses her in an oh-so-slightly imperceptible fashion. 

Jovannah, then tilting her head breaks the pregnancy: \textit{"Willard, you've been staring at my blouse for the last twenty seconds. If you're solving an isoperimetric problem in your head, I'd love to hear about it."}

Willard, jostling for his coffee as if they were in the dark: \textit{"I—what? I mean, well, it's just—the way fabric folds over a curved surface—it's, uh, quite like a minimal surface problem in differential geometry..."}

Jovannah, smirking: \textit{"Ah, so you're saying I’m an optimization problem or a problem to be solved?"}

Willard, thinking we are only a minute in and I am losing control; clearing his throat, scrambling for a footing: \textit{"I was merely thinking about how the tension distributes along the contours. It’s reminiscent of how spherical harmonics describe standing wave solutions on a sphere."}

Jovannah, leaning forward: \textit{"Spherical harmonics? Please, go on. I love it when you talk all eigenfunctions."}

Willard, not aware he is being led a merry dance feels he's on familiar ground now: \textit{"Well, the spherical harmonics $Y_{\ell m}$ are solutions to Laplace’s equation in spherical coordinates:"}
\begin{equation*}
\nabla^2 Y_{\ell m}(\theta, \phi) + \lambda Y_{\ell m}(\theta, \phi) = 0.
\end{equation*}
\textit{"They determine the natural vibrational modes of a sphere, much like how overtones define the timbre of a musical instrument. Just as the tension in a guitar string affects its harmonics so too the way your blouse drapes.}

% drumhead affects its harmonics, the way your blouse drapes—"}

Jovannah, raising an eyebrow: \textit{"So now my fashion choices are an eigenvalue problem?"}

Willard, back to being flustered: \textit{"No, no! That’s not what I meant. It’s just—take a vibrating drum. The way it resonates is governed by Bessel functions, much like spherical harmonics but in two dimensions."}

Jovannah, feigning deep thought: \textit{"So if my outfit were a drum, you'd be solving for the nodes of vibration?"}

Willard, now more thoroughly losing composure: \textit{"In a purely mathematical sense, yes, the modes of oscillation could be approximated - back of the nap-kinning - by Bessel functions, given by:"}
\begin{equation*}
 r^2 \frac{d^2 J_n}{dr^2} + r \frac{d J_n}{dr} + (r^2 - n^2) J_n = 0.
\end{equation*}
\textit{"These functions determine how waves propagate in circular domains, such as sound waves in a drum—or, well, other circular systems..."}

Jovannah, resting her chin on her hand, watching him squirm: \textit{"So, according to you, my blouse is a vibrating drum, my curves are dictated by eigenfunctions, and I’m basically a walking wave equation."}

Willard, realizing he's been backed into a corner: \textit{"That is... not entirely inaccurate?"}

Jovannah, laughing, placing a hand on his arm: \textit{"Relax, Willard. You’re adorable when you try to explain things without realizing where you're going with them. Pausing so he can misguidedly gather himself. Now, tell me ... if I were a perfect sphere, what would my fundamental mode of vibration be?"}

Willard, now finally getting the game and sighing in defeat, but taking on a formal tone despite himself: \textit{"$\ell = 0$, the monopole mode. Just a simple, uniform oscillation."}

Jovannah, winking, yes winking he was sure: \textit{"How disappointingly simple. Lucky for you, I have much more interesting oscillations."}
%%%%%%%%%%%%%
\bigskip

\textit{And then in came Ernest.}  Ernest Olber had a knack for appearing at precisely such moments. He strode in with his usual self-assurance, nodded to the waiter in a knowing manner for a coffee. As a man who had never met a mortal he couldn't lecture into submission every waiter he confronted was suitably responsive to his will.

Jovannah, never letting an opportunity pass, sidles towards Willard inviting Ernest to sit the other side, and a little unguarded for her, but grinning asks: \textit{"You know, you physicists always love a good harmonic. Go on Ernest, what’s so special about them?"}

\bigskip

Ernest, not needing such an invitation, sketches a quick diagram on a napkin: \textit{"Spherical harmonics describe standing waves on a sphere."}

\bigskip

Jovannah, feigning deep thought: \textit{"Harmonics... so the universe is just a very complicated steel drum?"}

\bigskip

Ernest, chuckling: \textit{"Not far off! The cosmic microwave background, ... the afterglow of the Big Bang, is actually analyzed in terms of spherical harmonics. We decompose its temperature fluctuations just like you decompose sound into different frequencies."}

\bigskip

Jovannah, tapping her fingers on the table: \textit{"Okay, so let’s take it a step further. What if we’re talking about something flat? Like ripples on a pond?"}

\bigskip

Ernest, nodding approvingly: \textit{"Yes. That’s where Bessel functions come in. They describe wave patterns in circular domains, like ripples in a drumhead. While spherical harmonics describe vibrations on a sphere, Bessel functions describe how waves behave in a circular boundary."}

\bigskip

Jovannah, leaning in while looking back in Willard's direction, voice playful: \textit{"So you're saying Ernest that my hemline behaves like a Bessel function?"}

\bigskip

Ernest, a little forlornly, seeing that he is being played: \textit{"Yes, Jovannah. Physics is everywhere. Even in your skirt"}

\bigskip

Jovannah, sitting back, satisfied, looking at Willard but addressing Ernest: \textit{"I caught Willard staring at my blouse earlier he claimed he was just solving an isoperimetric problem."}

Before Willard can summon a defence, Ernest, eager to reassert control offers derivatively: \textit{"Ah, Jovannah, what you’re really embodying is the isoperimetric inequality. It’s fundamental to physics. Soap bubbles form spheres because they minimize surface area for a given volume ..."}  

\bigskip

Jovannah, cutting him off, eyes glinting:  
\textit{"So my blouse clings the way it does because of… minimal surfaces?"}  

\bigskip

Willard, for once, finds his voice before Ernest can reclaim the floor.  \textit{"Not quite. It’s not just about minimization, ... it’s about harmonic distributions of tension. Which, actually, would bring us more pertinently to spherical harmonics."}  

\bigskip

Ernest, momentarily caught off guard, raises an eyebrow.  

\bigskip

Jovannah, feigning innocent curiosity to the extent that all parties know so: \textit{"Those spherical harmonics? Go on, Willard, you have my full attention."}

\bigskip

Willard, a little emboldened,  \textit{"Spherical harmonics $Y_{\ell m}$ describe standing wave patterns on a sphere so the way tension distributes in fabric over a curved surface is mathematically similar to how sound resonates in an enclosed space. As with a vibrating drum, but wrapped around a sphere."}  

\bigskip

Ernest, refusing to be outdone: \textit{"Ah, but if you’re talking about a drumhead, then you need Bessel functions, not spherical harmonics."}  

\bigskip

Jovannah, resting her chin on her hand, enjoying this immensely:  
\textit{"So according to you, Ernest, my hemline behaves like a Bessel function?"}

\bigskip

Ernest, uncharacteristically hesitating, sensing the trap:  \textit{"W-well, I...uh...yes? The way tension distributes itself along a curved hem could correspond to a standing wave pattern described by Bessel functions."}  

\bigskip

Jovannah turns back to Willard, her smile widening.  \textit{"And what do you think, Willard? Am I a drum, or a sphere?"}

\bigskip

Willard, regaining his composure as Ernest begins to wilt:  \textit{"Neither. You’re far too dynamic to be captured by a single eigenfunction... but if I had to choose, I’d say you’re an evolving system, wholly nonlinear, of higher order, beautifully unresolved."}  

\bigskip

Jovannah blinks, momentarily caught off guard.  

\bigskip

Ernest frowns. He' seen it, see the angular momentum of their system has  shifted.  

\bigskip

Jovannah, after a pause, eyes now a twinkling: \textit{"Willard, that might be the most charming thing anyone’s ever said to me."}  

\bigskip

Ernest, resigned that he has somehow lost this joust, relaxes, lips pressing into a thin line, sips his espresso that had been just dispatched by a waiter in knowing solemn silence.

\bigskip

\textit{(Somewhere, unseen, a Bessel function oscillates in quiet approval.)}

\bigskip

\begin{figure}[H]
    \centering
    \includegraphics[width=0.7\textwidth]{Illustrations/vign-whistfulWillard.png}
    \caption{Willard unspoken}
    \label{fig:pirallingprimes}
\end{figure}











































